% =====================================================================
% Wilson loops vs modular-projection surrogates; No-Go unchanged
% =====================================================================

\chapter{Wilson Loops, Surrogate Projections, and the No-Go}
\label{ch:wilson-surrogate}

\section{Wilson loops and the genuine area law}

Wilson loops are the fundamental gauge-invariant observables of lattice gauge
theory. On a discrete lattice one assigns a group element (the parallel
transporter) \(U_\ell\in G\) to every directed link \(\ell\). For a closed path
\(C\) the Wilson loop is
\begin{equation}
W(C)
=\operatorname{Tr}\,
\mathcal{P}
\prod_{\ell\in C} U_\ell.
\end{equation}
The \emph{area law} is the statement that the expectation value falls
exponentially with the area of the surface enclosed by \(C\):
\begin{equation}
\langle W(C)\rangle \sim e^{-\sigma\,\mathrm{Area}(C)}.
\end{equation}
The coefficient \(\sigma\) is the \emph{string tension}: a non-zero value means
the energy of a pair of static sources grows linearly with separation, with
\(\sigma\) the energy per unit length stored in the confining flux tube. A
\emph{perimeter law} \(\langle W\rangle\sim e^{-\mu\,\mathrm{Perimeter}(C)}\)
signals deconfinement or screening.

No genuine Wilson loop exists inside the resonant dynamics of the present model,
so a genuine area law and a genuine string tension are likewise absent: no gauge
flux tubes are defined.

\section{Absence of genuine Wilson loops in the resonant model}

In the resonant dynamics of the present model no continuum or discrete gauge
field is defined, and therefore no genuine Wilson loop exists. The modular
projections act on the integer state of the ternary map; they restore a residue
condition (\(Q=n\bmod 9\)) or a phase tolerance, but they do not assign
group-valued transporters to the edges of any lattice. Consequently the
sequences of projections cannot be assembled into a Wilson loop in the ordinary
sense.

\begin{theorem}[No Wilson loop from modular projections]
\label{thm:no-wilson}
Let \(\{n_t\}\) be a trajectory of \(T_3\) or \(T^\sharp\), and let projections
denote charge-preserving or min-defect corrections restoring \(Q=n\bmod 9\).
There does not exist a lattice gauge theory \((G,\{U_\ell\})\) constructed
solely from these data such that the ordered product of projections around a
closed residue sequence equals a Wilson loop \(W(C)\) in the sense of lattice
gauge theory.
\end{theorem}

\begin{proof}[Proof sketch]
A Wilson loop requires link variables in a group \(G\) on a spatial lattice.
The model supplies only integer states and residue/correction events on a
one-dimensional discrete-time orbit. No assignment \(\ell\mapsto U_\ell\in G\)
is given, so the trace of a path-ordered product of group elements is undefined.
\end{proof}

\section{Surrogate closed sequences}

\begin{definition}[Surrogate loop weight]
Regard each modular projection as a discrete correction event. A
\emph{surrogate closed sequence} is a closed walk in the residue graph
(e.g.\ on \(\mathbb{Z}/3\mathbb{Z}\) or \(\mathbb{Z}/27\mathbb{Z}\)) or a closed
pattern of projection-occurrence types along a trajectory. Let \(w(L,A)\) be a
statistical weight of such sequences classified by a perimeter proxy \(L\) and
an enclosed-step / area proxy \(A\).
\end{definition}

\begin{remark}[Area vs perimeter analogues; surrogate string tension]
If frequencies of closed sequences decay exponentially with the enclosed
step-count \(A\) (discrete ``area''), an area-law \emph{analogue} is present:
large closed patterns of projections or residue changes are statistically
suppressed. The decay constant \(\sigma_{\mathrm{surr}}\) of that exponential
may be called a \emph{surrogate string tension}; it measures only the strength
of that statistical suppression along constrained trajectories. If frequencies
decay only with the boundary length \(L\) (discrete ``perimeter''), a
perimeter-law analogue is present. Neither signal constitutes confinement of a
gauge charge or a genuine string tension, nor generates the integer \(539\) or
the period \(539.9\).
\end{remark}

\section{Downstream consequence, not source}

\begin{theorem}[Surrogate loops do not lift the No-Go]
\label{thm:surrogate-nogo}
Long trajectories on which surrogate loops are measured are produced only after
a fixed iteration count and a phase-locking period have already been imposed.
Any area-law or perimeter-law behaviour extracted from them is therefore a
consequence of those external constraints, not an independent dynamical origin
of the constraints. In particular it cannot serve as an independent source of
the integer \(539\) or of the period \(539.9\).

The a priori charge-correcting estimate for the completed map \(T^\sharp\)
continues to yield a short mean contraction depth of order \(N_\star=14\). The
No-Go Theorem continues to place the long resonant structure outside the reach
of residue structure, tower topology and flux democracy---whether or not
surrogate-loop weights are tallied along forced trajectories.

An area-law analogue (and any surrogate string tension extracted from it) is a
possible statistical signature of the constrained system and may be looked for
under the locked protocol. It remains a \emph{downstream consequence} (an
\emph{effect}) of the resonant apparatus, not a \emph{source} (a \emph{cause})
of that apparatus.
\end{theorem}

\begin{proof}
Immediate from Theorem~\ref{thm:nogo-canonical} (canonical No-Go) together with
the observation that surrogate weights are functionals of trajectories already
conditioned on external length / phase constraints when those are imposed.
\end{proof}

\section{Empirical permission}

\begin{corollary}[Empirical surrogate-loop study]
An empirical tally of surrogate-loop weights may characterise how modular
projections organise themselves along forced trajectories. It cannot lift the
No-Go or convert the resonant period into a quantity derived from the local map
alone.
\end{corollary}

% end
